\documentclass[a4paper]{article}
\usepackage[pdftex]{graphicx}
\usepackage[utf8]{inputenc}
\usepackage{enumerate}
\usepackage{icomma}
\usepackage{colortbl}
\usepackage{siunitx}
\sisetup{locale=DE}
\usepackage{amssymb}
\usepackage{tikz}
\usepackage{href-ul}
\hypersetup{
	colorlinks=true,
	linkcolor=blue,
	urlcolor=blue}
\usepackage{geometry}
\geometry{a4paper, top=15mm, left=15mm, right=15mm, bottom=15mm,
	headsep=10mm, footskip=12mm}

\begin{document}
	\begin{tikzpicture}(10,3)
		\draw[ultra thick](2,0) --(10,0) -- (10,1.3) --(2,1.3) -- (2,0);
		\draw[fill=black](2,0)-- (0,.6) -- (2,.6) -- (2,0);
		\node at (6,.6) {\large Solution sheet from \href{http://www.okuyakl.com}{www.okuyakl.com}};
	\end{tikzpicture}
	\vspace{0.5 cm}
	
	\noindent{\bf Task 1.}\\
	Both coordinates must satisfy the Pythagorean theorem, here: $x^2+y^2=1$.
	$$a)\, P(-1|0) \qquad b) \, P\left({1\over 2}|{1\over 2}\sqrt{3}\right) \qquad
	c)\, P\left({1 \over 2}\sqrt{2}|{1 \over 2}\sqrt{2}\right) \qquad d)\, P\left(-{1 \over 2 } \sqrt{3}|{1\over 2}\right)
	\qquad e) \, P(-0.6|0.8)$$
	
	\noindent{\bf Task 2.}\\
	$$a) \, \sin 30^\circ={1\over 2} \quad b)\, \cos 60^\circ= {1 \over 2} \quad c)\, \tan 45^\circ =1
	\quad d)\, \sin 45^\circ = {1\over 2} \sqrt{2} \quad e)\, \tan 60^\circ = \sqrt{3}$$
	
	\noindent{\bf Task 3.}\\
	The angle $\gamma$ is:
	$$\gamma=180^\circ-60^\circ-75^\circ=45^\circ$$
	This angle is divided into thirds, so every partial angle at point C is the same
	$\gamma_{1/2/3}=15^\circ$.
	In the first partial triangle from the left the missing angle is
	$\varepsilon_1= 180^\circ - 60^\circ -15^\circ = 105^\circ$.
	Using the law of sines we get here:
	$${c_1 \over \sin \gamma_1} = {b \over \sin \varepsilon_1} \qquad \Leftrightarrow
	\qquad c_1 = {b \over \sin \varepsilon_1} \cdot \sin \gamma_1 = {8\, \SI{}{\centi \meter} \over \sin 105^\circ} \cdot
	\sin 15^\circ = 2.14 \, \SI{}{\centi\meter}$$
	The area of this first partial triangle is:
	$$A_1={1\over 2} \cdot c_1 \cdot b \cdot \sin \alpha = 0.5 \cdot 2.14 \cdot 8 \cdot \sin 60^\circ =
	7.41\,\SI{}{\centi\meter}^3$$
	We need the side $x$ for further calculations; we get them using the law of sines:
	$${x \over \sin \alpha} = {b \over \sin \varepsilon_1} \qquad \Leftrightarrow
	\qquad x = {b \over \sin \varepsilon_1} \cdot \sin \alpha = {8\, \SI{}{\centi \meter} \over \sin 105^\circ} \cdot
	\sin 60^\circ = 7.17 \, \SI{}{\centi\meter}$$
	In the second partial triangle the missing angle is
	$\varepsilon_2= 180^\circ -15^\circ-(180^\circ-105^\circ)=90^\circ$.
	So the triangle is right-angled.
	With the sine we get here:
	$${c_2 \over x}= \sin \gamma_2 \qquad \Leftrightarrow
	\qquad c_2 = x \cdot \sin \gamma_2 = 7.17\, \SI{}{\centi\meter} \cdot \sin 15^\circ = 1.86 \, \SI{}{\centi\meter }$$
	The area of this second partial triangle is:
	$$A_2={1\over 2} \cdot c_2 \cdot x \cdot \sin (180^\circ-105^\circ) = 0.5 \cdot 1.86 \cdot 7.17 \cdot \sin 75 ^\circ =6.44\,\SI{}{\centi\meter}^3$$
	According to Pythagoras, the side $y$ is:
	$$y^2=x^2-c_2^2 \qquad \Leftrightarrow \qquad y=\sqrt{x^2-c_2^2-}=\sqrt{47.95\,\SI{}{\centi\ meter}^2}=6.92\,\SI{}{\centi\meter}$$
	We get $A_3$ by considering that the middle and right triangles are congruent,
	Together they form an isosceles triangle with height $y$ and base $c_2 + c_3 =2 c_2$
	so $$A_3=A_2=6.44\,\SI{}{\centi\meter}^3$$
	
	\renewcommand{\arraystretch}{3}
	
	\begin{tabular}{|>{\columncolor[gray]{.8}}p{2 cm}|p{2 cm}|p{2 cm}|p{2 cm}|p{2 cm}|p {2cm}|p{2cm}|}
		\hline
		\rowcolor[gray]{.8}
		& \bf {a}&\bf b&\bf c& $ \alpha$ & $ \beta$ &$ \gamma$\\
		\hline
		a)& 7 & 5 & 5.8 &$80^\circ$ & $45^\circ$ & $55^\circ$\\
		\hline
		b)& 3,6 & 8&6 &$25^\circ$ & $111^\circ$ & $44^\circ$\\
		\hline
		c)& 10&10,1 & 7,2 & $68^\circ$& $70^\circ$& $42^\circ$ \\
		\hline
		d)& 6.1 &5.3 &6.5 &$61^\circ$ &$50^\circ$ & $69^\circ$\\
		\hline
		e)& 1,1 & 12 &11,8 &$5^\circ$ & $100^\circ$ & $75^\circ$\\
		\hline
		f)& 4.8&2.9 &3.5 &$98^\circ$ &$36^\circ$ &$46^\circ$ \\
		\hline
		g)& 4.8 & 9.1 & & $103^\circ$ & &\\
		\hline
		h)& 14 & & &$95^\circ$ & $100^\circ$ &\\
		\hline
	\end{tabular}
	
	\begin{center}
		\includegraphics[width=7 cm]{../../viecher/eendcomic.pdf}
		
		Here you can go back to the \href{https://www.okuyakl.de/math/m10trigL035/ae035.pdf}{task sheet}
	\end{center}
	
\end{document}